\section*{Chapitre \ref{ch:qm:ito-calculus} --- Calcul d'Itô}

\begin{solution}{exo:qm:ito-calculus:1}
$\E\int_0^1W_t^2\,dt = \int_0^1t\,dt = \tfrac12$. Grâce à l'\cref{ex:qm:ito-calculus:wdw},
$\Var(\tfrac12(W_1^2 - 1)) = \tfrac14\Var(W_1^2) = \tfrac14 \times 2 = \tfrac12$, et la moyenne est nulle.
\end{solution}

\begin{solution}{exo:qm:ito-calculus:2}
$d(W^3) = 3W^2\,dW + 3W\,dt$. En prenant les espérances, $\E[W_t^3] = 3\int_0^t\E[W_s]\,ds = 0$, et
$W_t^3 - 3\int_0^tW_s\,ds$ est une \omterm{def:qm:probability-at-speed:martingale}{martingale}.
\end{solution}

\begin{solution}{exo:qm:ito-calculus:3}
Moyenne $e^{1} = 2.72$, médiane $e^{(0.10 - 0.08) \times 10} = 1.22$. La médiane est constante lorsque $\sigma^2 = 2\mu$,
$\sigma = \sqrt{0.2} = 44.7\%$.
\end{solution}

\begin{solution}{exo:qm:ito-calculus:4}
C'est $\mathcal E(W)_t$~: $dZ = Z\,dW$, et $\E\int_0^tZ_s^2\,ds = \int_0^te^s\,ds < \infty$, donc l'\omterm{def:qm:ito-calculus:integral}{intégrale d'Itô}
est une véritable \omterm{def:qm:probability-at-speed:martingale}{martingale}. $\E[e^{2W_1}] = e^{2}= 7.39$ (la fonction génératrice des moments de $\mathcal
N(0,1)$ en 2).
\end{solution}

\begin{solution}{exo:qm:ito-calculus:5}
$d(tW_t) = W_t\,dt + t\,dW_t$, donc $\int_0^1W_t\,dt = W_1 - \int_0^1t\,dW_t = \int_0^1(1 - t)\,dW_t$, de variance
$\int_0^1(1 - t)^2dt = \tfrac13$ par l'isométrie.
\end{solution}

\begin{solution}{exo:qm:ito-calculus:6}
$[X, Y]_t = 2t$, $[Y]_t = 5t$, corrélation $2/\sqrt5 = 0.894$.
\end{solution}

\begin{solution}{exo:qm:ito-calculus:7}
$L = 60$ donne $a = 0.223$ et un ratio de Sharpe avec look-ahead de 3.45~; la moyenne sur 100
marches aléatoires initialisées par une graine est de 3.47. Une moyenne plus lente accorde moins de poids au dernier accroissement, donc le biais est plus faible,
mais il reste énorme.
\end{solution}

\begin{solution}{exo:qm:ito-calculus:8}
Le signal utilise le cours de clôture, qui n'est pas connu avant l'enchère de clôture au cours de laquelle la
règle prétend négocier a été publié~: la décision n'est pas adaptée à l'information disponible
au moment où l'ordre doit être envoyé. Calculer le signal à partir du cours indicatif de l'enchère publié
avant la limite de transmission de l'ordre (chapitre~1), ou négocier à la prochaine ouverture.
\end{solution}

\begin{solution}{pb:qm:ito-calculus:1}
\textbf{1.} $c = 0.01\,(\sum_{m=1}^{19}(m/20)^2)^{1/2} = 0.01 \times 2.485 = 0.0248$.
\textbf{2.} $a = 0.95/2.485 = 0.382$.
\textbf{3.} Moyenne $a\sigma = 0.0038$ et écart-type $\sigma\sqrt{1 + a^2} = 0.0107$ par jour, par
unité de position.
\textbf{4.} $\sqrt{252} \times 0.382/1.070 = 5.67$.
\textbf{5.} Zéro~: $\theta_{t}$ est indépendant de $S_{t+1} - S_t$.
\textbf{6.} 5.43 même barre, $-0.04$ barre suivante.
\textbf{7.} Même barre $5.69 \pm 0.13$~; barre suivante $0.00 \pm 0.32$.
\textbf{8.} Les deux valent 9.653 (965~\% du notionnel)~; la différence est $2 \times 10^{-14}$.
\textbf{9.} 34.9.
\textbf{10.} Sans edge, le ratio de Sharpe annualisé estimé sur $T$ années a une erreur type
d'environ $1/\sqrt T$ (chapitre~11)~: $0.32$ pour dix ans.
\textbf{11.} $\sum\theta_t\Delta S_t = \sum\theta_{t-1}\Delta S_t + \sum\Delta\theta_t\Delta S_t$~: une somme d'Itô plus la
covariation $[\theta, S]$.
\textbf{12.} Son intégrande est adapté, c'est donc une \omterm{def:qm:probability-at-speed:martingale}{martingale} partant de zéro.
\textbf{13.} 9.36, 7.47 et 3.45.
\textbf{14.} Le ratio de Sharpe quotidien par barre est inchangé, et l'annualisation sur $252 \times 78$ barres
le multiplie par $\sqrt{78}$~: 50.
\textbf{15.} Oui, si le signal n'utilise que le cours de clôture et que les gains commencent à l'ouverture suivante~; la
position manque alors le mouvement de nuit, que le backtest honnête doit également exclure.
\textbf{16.} Décaler les positions d'une barre et vérifier si la performance s'effondre~; lancer
\texttt{lookahead\_test}, ou régresser le P\&L quotidien sur $\Delta\theta_t\Delta S_t$~; et lire les
horodatages des entrées du signal.
\textbf{17.} Non~: son signe est celui de $[\theta, S]$~; une règle de retour à la moyenne, qui vend lorsque le
prix monte, présente un biais de look-ahead négatif.
\textbf{18.} Le rebond bid-ask, entre cours acheteur et cours vendeur, fait revenir vers la moyenne les prix de dernière transaction~; une règle honnête
de retour à la moyenne sur ces prix «~gagne~» le rebond, que l'on ne capte qu'en payant la fourchette
(le livre~10).
\textbf{19.} Résultat nommé~: \emph{le biais de look-ahead dans le backtest}~: sur une marche aléatoire, le backtest en même barre
d'une règle à variance unitaire corrélée à hauteur de $a$ au dernier accroissement, a un ratio de Sharpe quotidien
égal à $a/\sqrt{1 + a^2}$~; pour la règle de tendance sur 20 jours, $a = 0.382$ et le ratio de Sharpe annualisé vaut 5.67
(5.69 en simulation), contre zéro pour la somme d'Itô.
\textbf{20.} Chaque position est fixée à partir de l'information disponible avant le mouvement de prix
par lequel elle est multipliée.
\end{solution}

\begin{solution}{iq:qm:ito-calculus:1}
$\tfrac12(W_T^2 - T)$. Les sommes au point gauche diffèrent de $\tfrac12\sum\Delta(W^2)$ d'une quantité $\tfrac12\sum(\Delta W)^2$,
qui tend vers $T$~: la \omterm{def:qm:brownian-motion:qv}{variation quadratique} ne s'annule pas.
\iqlookfor{l'identité télescopique et la \omterm{def:qm:brownian-motion:qv}{variation quadratique}.}
\end{solution}

\begin{solution}{iq:qm:ito-calculus:2}
$\mu - \sigma^2/2 = 2\%$ par an. La moyenne de $S_{10}/S_0$ vaut $e^{1} = 2.72$ mais la médiane vaut $1.22$, et un
dixième des détenteurs finit sous 0.24~: la moyenne est portée par quelques issues de grande ampleur.
\iqlookfor{la correction d'Itô et la moyenne par rapport à la médiane.}
\end{solution}

\begin{solution}{iq:qm:ito-calculus:3}
Non~: $d(W^2) = 2W\,dW + dt$, donc $W_t^2 - t$ est la \omterm{def:qm:probability-at-speed:martingale}{martingale}. $\Var(\int_0^TW\,dW) = \int_0^Tt\,dt =
T^2/2$.
\iqlookfor{la formule d'Itô et l'isométrie.}
\end{solution}

\begin{solution}{iq:qm:ito-calculus:4}
Piloter le backtest par événements, dans l'ordre temporel~: chaque donnée porte l'instant où elle est devenue disponible, la
stratégie voit une vue as-of, et un ordre émis en $t$ ne peut être exécuté que contre des prix postérieurs à
$t$ plus une latence. Le moteur, et non la stratégie, applique les positions à l'intervalle suivant, et
un test compare chaque simulation à une simulation décalée et signale la covariation entre les positions et
rendements contemporains.
\iqlookfor{horodatages de disponibilité, jointures as-of et décalage structurel.}
\end{solution}

\begin{solution}{iq:qm:ito-calculus:5}
Un processus qui n'est une \omterm{def:qm:probability-at-speed:martingale}{martingale} que jusqu'à une suite de \omterm{def:qm:probability-at-speed:stopping}{temps d'arrêt}~; un processus positif est une
\omterm{def:qm:probability-at-speed:martingale}{surmartingale} et peut perdre son espérance, comme $1/|B_t|$ pour un \omterm{def:qm:brownian-motion:bm}{mouvement brownien} tridimensionnel.
Une densité candidate pour un \omterm{def:qm:probability-at-speed:change}{changement de probabilité}, ou un prix déflaté, qui est une \omterm{def:qm:ito-calculus:local}{martingale locale}
stricte donne des prix qui violent la parité call-put ou une mesure qui n'est pas une probabilité.
\iqlookfor{localisation, lemme de Fatou et conséquence pour la valorisation.}
\end{solution}

\begin{solution}{iq:qm:ito-calculus:6}
La finance utilise la version d'Itô~: une position doit être fixée avant le mouvement de prix, qui est le point gauche, et
les gains qui en résultent sont des \omterm{def:qm:probability-at-speed:martingale}{martingales} sous la bonne probabilité. La physique utilise souvent Stratonovich,
la limite d'un bruit régulier, car elle conserve la règle de la chaîne ordinaire. Elles diffèrent d'un terme $\tfrac12[X,W]$.
\iqlookfor{non-anticipation face aux limites de bruit régulier.}
\end{solution}
